% Options for packages loaded elsewhere
% Options for packages loaded elsewhere
\PassOptionsToPackage{unicode}{hyperref}
\PassOptionsToPackage{hyphens}{url}
\PassOptionsToPackage{dvipsnames,svgnames,x11names}{xcolor}
%
\documentclass[
  letterpaper,
  DIV=11,
  numbers=noendperiod]{scrartcl}
\usepackage{xcolor}
\usepackage{amsmath,amssymb}
\setcounter{secnumdepth}{-\maxdimen} % remove section numbering
\usepackage{iftex}
\ifPDFTeX
  \usepackage[T1]{fontenc}
  \usepackage[utf8]{inputenc}
  \usepackage{textcomp} % provide euro and other symbols
\else % if luatex or xetex
  \usepackage{unicode-math} % this also loads fontspec
  \defaultfontfeatures{Scale=MatchLowercase}
  \defaultfontfeatures[\rmfamily]{Ligatures=TeX,Scale=1}
\fi
\usepackage{lmodern}
\ifPDFTeX\else
  % xetex/luatex font selection
\fi
% Use upquote if available, for straight quotes in verbatim environments
\IfFileExists{upquote.sty}{\usepackage{upquote}}{}
\IfFileExists{microtype.sty}{% use microtype if available
  \usepackage[]{microtype}
  \UseMicrotypeSet[protrusion]{basicmath} % disable protrusion for tt fonts
}{}
\makeatletter
\@ifundefined{KOMAClassName}{% if non-KOMA class
  \IfFileExists{parskip.sty}{%
    \usepackage{parskip}
  }{% else
    \setlength{\parindent}{0pt}
    \setlength{\parskip}{6pt plus 2pt minus 1pt}}
}{% if KOMA class
  \KOMAoptions{parskip=half}}
\makeatother
% Make \paragraph and \subparagraph free-standing
\makeatletter
\ifx\paragraph\undefined\else
  \let\oldparagraph\paragraph
  \renewcommand{\paragraph}{
    \@ifstar
      \xxxParagraphStar
      \xxxParagraphNoStar
  }
  \newcommand{\xxxParagraphStar}[1]{\oldparagraph*{#1}\mbox{}}
  \newcommand{\xxxParagraphNoStar}[1]{\oldparagraph{#1}\mbox{}}
\fi
\ifx\subparagraph\undefined\else
  \let\oldsubparagraph\subparagraph
  \renewcommand{\subparagraph}{
    \@ifstar
      \xxxSubParagraphStar
      \xxxSubParagraphNoStar
  }
  \newcommand{\xxxSubParagraphStar}[1]{\oldsubparagraph*{#1}\mbox{}}
  \newcommand{\xxxSubParagraphNoStar}[1]{\oldsubparagraph{#1}\mbox{}}
\fi
\makeatother

\usepackage{color}
\usepackage{fancyvrb}
\newcommand{\VerbBar}{|}
\newcommand{\VERB}{\Verb[commandchars=\\\{\}]}
\DefineVerbatimEnvironment{Highlighting}{Verbatim}{commandchars=\\\{\}}
% Add ',fontsize=\small' for more characters per line
\usepackage{framed}
\definecolor{shadecolor}{RGB}{241,243,245}
\newenvironment{Shaded}{\begin{snugshade}}{\end{snugshade}}
\newcommand{\AlertTok}[1]{\textcolor[rgb]{0.68,0.00,0.00}{#1}}
\newcommand{\AnnotationTok}[1]{\textcolor[rgb]{0.37,0.37,0.37}{#1}}
\newcommand{\AttributeTok}[1]{\textcolor[rgb]{0.40,0.45,0.13}{#1}}
\newcommand{\BaseNTok}[1]{\textcolor[rgb]{0.68,0.00,0.00}{#1}}
\newcommand{\BuiltInTok}[1]{\textcolor[rgb]{0.00,0.23,0.31}{#1}}
\newcommand{\CharTok}[1]{\textcolor[rgb]{0.13,0.47,0.30}{#1}}
\newcommand{\CommentTok}[1]{\textcolor[rgb]{0.37,0.37,0.37}{#1}}
\newcommand{\CommentVarTok}[1]{\textcolor[rgb]{0.37,0.37,0.37}{\textit{#1}}}
\newcommand{\ConstantTok}[1]{\textcolor[rgb]{0.56,0.35,0.01}{#1}}
\newcommand{\ControlFlowTok}[1]{\textcolor[rgb]{0.00,0.23,0.31}{\textbf{#1}}}
\newcommand{\DataTypeTok}[1]{\textcolor[rgb]{0.68,0.00,0.00}{#1}}
\newcommand{\DecValTok}[1]{\textcolor[rgb]{0.68,0.00,0.00}{#1}}
\newcommand{\DocumentationTok}[1]{\textcolor[rgb]{0.37,0.37,0.37}{\textit{#1}}}
\newcommand{\ErrorTok}[1]{\textcolor[rgb]{0.68,0.00,0.00}{#1}}
\newcommand{\ExtensionTok}[1]{\textcolor[rgb]{0.00,0.23,0.31}{#1}}
\newcommand{\FloatTok}[1]{\textcolor[rgb]{0.68,0.00,0.00}{#1}}
\newcommand{\FunctionTok}[1]{\textcolor[rgb]{0.28,0.35,0.67}{#1}}
\newcommand{\ImportTok}[1]{\textcolor[rgb]{0.00,0.46,0.62}{#1}}
\newcommand{\InformationTok}[1]{\textcolor[rgb]{0.37,0.37,0.37}{#1}}
\newcommand{\KeywordTok}[1]{\textcolor[rgb]{0.00,0.23,0.31}{\textbf{#1}}}
\newcommand{\NormalTok}[1]{\textcolor[rgb]{0.00,0.23,0.31}{#1}}
\newcommand{\OperatorTok}[1]{\textcolor[rgb]{0.37,0.37,0.37}{#1}}
\newcommand{\OtherTok}[1]{\textcolor[rgb]{0.00,0.23,0.31}{#1}}
\newcommand{\PreprocessorTok}[1]{\textcolor[rgb]{0.68,0.00,0.00}{#1}}
\newcommand{\RegionMarkerTok}[1]{\textcolor[rgb]{0.00,0.23,0.31}{#1}}
\newcommand{\SpecialCharTok}[1]{\textcolor[rgb]{0.37,0.37,0.37}{#1}}
\newcommand{\SpecialStringTok}[1]{\textcolor[rgb]{0.13,0.47,0.30}{#1}}
\newcommand{\StringTok}[1]{\textcolor[rgb]{0.13,0.47,0.30}{#1}}
\newcommand{\VariableTok}[1]{\textcolor[rgb]{0.07,0.07,0.07}{#1}}
\newcommand{\VerbatimStringTok}[1]{\textcolor[rgb]{0.13,0.47,0.30}{#1}}
\newcommand{\WarningTok}[1]{\textcolor[rgb]{0.37,0.37,0.37}{\textit{#1}}}

\usepackage{longtable,booktabs,array}
\newcounter{none} % for unnumbered tables
\usepackage{calc} % for calculating minipage widths
% Correct order of tables after \paragraph or \subparagraph
\usepackage{etoolbox}
\makeatletter
\patchcmd\longtable{\par}{\if@noskipsec\mbox{}\fi\par}{}{}
\makeatother
% Allow footnotes in longtable head/foot
\IfFileExists{footnotehyper.sty}{\usepackage{footnotehyper}}{\usepackage{footnote}}
\makesavenoteenv{longtable}
\usepackage{graphicx}
\makeatletter
\newsavebox\pandoc@box
\newcommand*\pandocbounded[1]{% scales image to fit in text height/width
  \sbox\pandoc@box{#1}%
  \Gscale@div\@tempa{\textheight}{\dimexpr\ht\pandoc@box+\dp\pandoc@box\relax}%
  \Gscale@div\@tempb{\linewidth}{\wd\pandoc@box}%
  \ifdim\@tempb\p@<\@tempa\p@\let\@tempa\@tempb\fi% select the smaller of both
  \ifdim\@tempa\p@<\p@\scalebox{\@tempa}{\usebox\pandoc@box}%
  \else\usebox{\pandoc@box}%
  \fi%
}
% Set default figure placement to htbp
\def\fps@figure{htbp}
\makeatother





\setlength{\emergencystretch}{3em} % prevent overfull lines

\providecommand{\tightlist}{%
  \setlength{\itemsep}{0pt}\setlength{\parskip}{0pt}}



 
\usepackage[]{biblatex}


\KOMAoption{captions}{tableheading}
\makeatletter
\@ifpackageloaded{caption}{}{\usepackage{caption}}
\AtBeginDocument{%
\ifdefined\contentsname
  \renewcommand*\contentsname{Table of contents}
\else
  \newcommand\contentsname{Table of contents}
\fi
\ifdefined\listfigurename
  \renewcommand*\listfigurename{List of Figures}
\else
  \newcommand\listfigurename{List of Figures}
\fi
\ifdefined\listtablename
  \renewcommand*\listtablename{List of Tables}
\else
  \newcommand\listtablename{List of Tables}
\fi
\ifdefined\figurename
  \renewcommand*\figurename{Figure}
\else
  \newcommand\figurename{Figure}
\fi
\ifdefined\tablename
  \renewcommand*\tablename{Table}
\else
  \newcommand\tablename{Table}
\fi
}
\@ifpackageloaded{float}{}{\usepackage{float}}
\floatstyle{ruled}
\@ifundefined{c@chapter}{\newfloat{codelisting}{h}{lop}}{\newfloat{codelisting}{h}{lop}[chapter]}
\floatname{codelisting}{Listing}
\newcommand*\listoflistings{\listof{codelisting}{List of Listings}}
\makeatother
\makeatletter
\makeatother
\makeatletter
\@ifpackageloaded{caption}{}{\usepackage{caption}}
\@ifpackageloaded{subcaption}{}{\usepackage{subcaption}}
\makeatother
\makeatletter
\@ifpackageloaded{fontawesome5}{}{\usepackage{fontawesome5}}
\makeatother
\usepackage{bookmark}
\IfFileExists{xurl.sty}{\usepackage{xurl}}{} % add URL line breaks if available
\urlstyle{same}
% fallback for those not using the hyperref driver hyperxmp:
\makeatletter
\@ifundefined{xmpquote}{\newcommand{\xmpquote}[1]{#1}}{}
\makeatother
\hypersetup{
  pdftitle={Pricing Derivatives Using Black-Scholes-Merton Model},
  pdfauthor={Rafiq Islam},
  colorlinks=true,
  linkcolor={blue},
  filecolor={Maroon},
  citecolor={Blue},
  urlcolor={Blue},
  pdfcreator={LaTeX via pandoc}}


\title{Pricing Derivatives Using Black-Scholes-Merton Model}
\author{Rafiq Islam}
\date{2023-11-12}
\begin{document}
\maketitle

\renewcommand*\contentsname{Table of contents}
{
\setcounter{tocdepth}{4}
\tableofcontents
}

\subsection{Introduction}\label{introduction}

In this blog, we will explore how to price simple equity derivatives
using the Black-Scholes-Merton (BSM) model. We will derive the
mathematical formula and then provide Python code to implement it.

\subsubsection{Background and
Preliminaries}\label{background-and-preliminaries}

Before proceeding to the deep of the discussion, we need to know some
definition and terminology

\textbf{Brownian Motion:} Brownian motion is a concept with definitions
and applications across various disciplines, named after the botanist
Robert Brown, is the random, erratic movement of particles suspended in
a fluid (liquid or gas) due to their collisions with the fast-moving
molecules of the fluid.

\emph{Brownian motion is a stochastic process \((B_t)_{t \geq 0}\)
defined as a continuous-time process with the following properties:}

\begin{itemize}
\tightlist
\item
  \(B_0 = 0\) almost surely.
\item
  \(B_t\) has independent increments.
\item
  For \(t > s\), \(B_t - B_s \sim N(0, t-s)\) (normally distributed with
  mean 0 and variance \(t-s\)).
\item
  \(B_t\) has continuous paths almost surely.
\end{itemize}

\begin{Shaded}
\begin{Highlighting}[]
\ImportTok{from}\NormalTok{ mywebstyle }\ImportTok{import}\NormalTok{ plot\_style}
\NormalTok{plot\_style(}\StringTok{\textquotesingle{}\#f4f4f4\textquotesingle{}}\NormalTok{)}
\ImportTok{import}\NormalTok{ numpy }\ImportTok{as}\NormalTok{ np}
\ImportTok{import}\NormalTok{ matplotlib.pyplot }\ImportTok{as}\NormalTok{ plt}

\CommentTok{\# Parameters}
\NormalTok{n\_steps }\OperatorTok{=} \DecValTok{100}  \CommentTok{\# Number of steps}
\NormalTok{n\_paths }\OperatorTok{=} \DecValTok{20}   \CommentTok{\# Number of paths}
\NormalTok{time\_horizon }\OperatorTok{=} \DecValTok{1}  \CommentTok{\# Total time}
\NormalTok{dt }\OperatorTok{=}\NormalTok{ time\_horizon }\OperatorTok{/}\NormalTok{ n\_steps  }\CommentTok{\# Time step}
\NormalTok{t }\OperatorTok{=}\NormalTok{ np.linspace(}\DecValTok{0}\NormalTok{, time\_horizon, n\_steps)  }\CommentTok{\# Time array}

\CommentTok{\# Generate Brownian motion}
\KeywordTok{def}\NormalTok{ generate\_brownian\_paths(n\_paths, n\_steps, dt):}
    \CommentTok{\# Standard normal increments scaled by sqrt(dt)}
\NormalTok{    increments }\OperatorTok{=}\NormalTok{ np.random.normal(}\DecValTok{0}\NormalTok{, np.sqrt(dt), (n\_paths, n\_steps))}
    \CommentTok{\# Cumulative sum to generate paths}
    \ControlFlowTok{return}\NormalTok{ np.cumsum(increments, axis}\OperatorTok{=}\DecValTok{1}\NormalTok{)}

\CommentTok{\# Generate one path and multiple paths}
\NormalTok{single\_path }\OperatorTok{=}\NormalTok{ generate\_brownian\_paths(}\DecValTok{1}\NormalTok{, n\_steps, dt)[}\DecValTok{0}\NormalTok{]}
\NormalTok{multiple\_paths }\OperatorTok{=}\NormalTok{ generate\_brownian\_paths(n\_paths, n\_steps, dt)}

\CommentTok{\# Plotting}
\NormalTok{fig, axes }\OperatorTok{=}\NormalTok{ plt.subplots(}\DecValTok{1}\NormalTok{, }\DecValTok{2}\NormalTok{, figsize}\OperatorTok{=}\NormalTok{(}\FloatTok{7.9}\NormalTok{, }\FloatTok{3.9}\NormalTok{))}

\CommentTok{\# Single path}
\NormalTok{axes[}\DecValTok{0}\NormalTok{].plot(t, single\_path, label}\OperatorTok{=}\StringTok{"Single Path"}\NormalTok{)}
\NormalTok{axes[}\DecValTok{0}\NormalTok{].set\_title(}\StringTok{"Brownian Motion: Single Path"}\NormalTok{)}
\NormalTok{axes[}\DecValTok{0}\NormalTok{].set\_xlabel(}\StringTok{"Time"}\NormalTok{)}
\NormalTok{axes[}\DecValTok{0}\NormalTok{].set\_ylabel(}\StringTok{"Position"}\NormalTok{)}
\NormalTok{axes[}\DecValTok{0}\NormalTok{].legend()}

\CommentTok{\# Multiple paths}
\ControlFlowTok{for}\NormalTok{ path }\KeywordTok{in}\NormalTok{ multiple\_paths:}
\NormalTok{    axes[}\DecValTok{1}\NormalTok{].plot(t, path, alpha}\OperatorTok{=}\FloatTok{0.5}\NormalTok{, linewidth}\OperatorTok{=}\FloatTok{0.8}\NormalTok{)}
\NormalTok{axes[}\DecValTok{1}\NormalTok{].set\_title(}\SpecialStringTok{f"Brownian Motion: }\SpecialCharTok{\{}\NormalTok{n\_paths}\SpecialCharTok{\}}\SpecialStringTok{ Paths"}\NormalTok{)}
\NormalTok{axes[}\DecValTok{1}\NormalTok{].set\_xlabel(}\StringTok{"Time"}\NormalTok{)}
\NormalTok{axes[}\DecValTok{1}\NormalTok{].set\_ylabel(}\StringTok{"Position"}\NormalTok{)}

\NormalTok{plt.tight\_layout()}
\NormalTok{plt.show()}
\end{Highlighting}
\end{Shaded}

\pandocbounded{\includegraphics[keepaspectratio]{index_files/figure-pdf/cell-2-output-1.pdf}}

\textbf{Geometric Brownian Motion (GBM)}\\
A stochastic process \(S_t\) is said to follow a geometric Brownian
motion if it satisfies the following equation:\\
\[
dS_t = \mu S_t dt+\sigma S_t dB_t
\]

Which can be written as\\
\[
S_t - S_0 =\int_0^t \mu S_u du + \int_0^t \sigma S_u dB_u
\]

To solve the GBM, we apply Ito's formula to the function
\(Z_t = f(t, S_t)= \ln(S_t)\) and then by Taylor's expansion, we have

\begin{align*}
df & = \frac{\partial f}{\partial t}dt+ \frac{\partial f}{\partial s}dS_t + \frac{1}{2} \frac{\partial ^2f}{\partial s^2}(dS_t)^2+\frac{1}{2}\frac{\partial ^2f}{\partial s^2}(dt)^2+\frac{\partial^2 f}{\partial t\partial s}dtdS_t
\end{align*}

By definition we have \begin{align*}
dS_t &= \mu S_t dt+\sigma S_t dB_t\\
(dS_t)^2 & = \mu^2 (dt)^2+2\mu \sigma dt dB_t + \sigma^2 (dB_t)^2
\end{align*}

The term \((dt)^2\) is negligible compared to the term \(dt\) and it is
also assume that the product \(dtdB_t\) is negligible. Furthermore, the
quadratic variation of \(B_t\) i.e., \((dB_t)^2= dt\). With these
values, we obtain

\begin{align*}
dZ_t & = \frac{1}{S_t} dS_t + \frac{1}{2} \left\{-\frac{1}{S_t^2}\right\}[dS_t]^2\\
& =  \frac{1}{S_t} (\mu S_t dt+\sigma S_t dB_t) + \frac{1}{2} \left\{-\frac{1}{S_t^2}\right\}\sigma^2S_t^2dt\\
\implies dZ_t &= (\mu dt +\sigma dB_t) -\frac{1}{2}\sigma^2 dt\\
& = \left(\mu-\frac{1}{2}\sigma^2\right)dt+\sigma dB_t
\end{align*}

with \(Z_0=\ln S_0\). Now we have the following\\
\begin{align*}
\int_0^t dZ_s &= \int_0^t \left(\mu-\frac{1}{2}\sigma^2\right)ds + \int_0^t\sigma dB_s\\
\implies Z_t - Z_0 &= \left(\mu-\frac{1}{2}\sigma^2\right)t + \sigma B_t\\
\implies \ln S_t - \ln S_0&= \left(\mu-\frac{1}{2}\sigma^2\right)t + \sigma B_t\\
\implies \ln \left(\frac{S_t}{S_0}\right) &= \left(\mu-\frac{1}{2}\sigma^2\right)t + \sigma B_t\\
\implies S_t &= S_0 \exp{\left\{\left(\mu-\frac{1}{2}\sigma^2\right)t + \sigma B_t\right\}}
\end{align*}

\subsection{Black-Scholes-Merton
Formula}\label{black-scholes-merton-formula}

Now we are ready to derive the BSM PDE. The payoff of an \emph{option}
\(V(S,T)\) at maturity is is known. To find the value at an earlier
stage, we need to know how V behaves as a function of \(S\) and \(t\).
By Ito's lemma we have \begin{align*}
dV& = \left(\mu S\frac{\partial V}{\partial S}+\frac{\partial V}{\partial t}+\frac{1}{2}\sigma^2S^2\frac{\partial^2 V}{\partial S^2}\right)dt+\sigma S\frac{\partial V}{\partial S}dB.
\end{align*}

Now let's consider a portfolio consisting of a short one option and long
\(\frac{\partial V}{\partial S}\) shares at time \(t\). The value of
this portfolio is

\[
\Pi = -V+\frac{\partial V}{\partial S}S
\]

over the time \([t,t+\Delta t]\), the total profit or loss from the
changes in the values of the portfolio is \[
\Delta \Pi = -\Delta V + \frac{\partial V}{\partial S}\Delta S
\]

Now by the discretization we have, \begin{align*}
\Delta S &= \mu S \Delta t +\sigma S \Delta B\\
\Delta V & = \left(\mu S\frac{\partial V}{\partial S}+\frac{\partial V}{\partial t}+\frac{1}{2}\sigma^2S^2\frac{\partial^2 V}{\partial S^2}\right)\Delta t+\sigma S\frac{\partial V}{\partial S}\Delta B\\
\implies \Delta \Pi  & = \left(-\frac{\partial V}{\partial t} -\frac{1}{2}\sigma^2S^2\frac{\partial^2 V}{\partial S^2}\right)\Delta t
\end{align*}

At this point, if \(r\) is the risk-free interest rate then we will have
following relationship \[
r\Pi \Delta t = \Delta \Pi 
\]

The rationale of this relation is that no-aribtrage assumption. Thus, we
have\\
\begin{align*}
\left(-\frac{\partial V}{\partial t} -\frac{1}{2}\sigma^2S^2\frac{\partial^2 V}{\partial S^2}\right)\Delta t & = r \left(- V + \frac{\partial V}{\partial S} S\right)\Delta t \\
\implies \frac{\partial V}{\partial t} + \frac{1}{2} \sigma^2 S^2 \frac{\partial^2 V}{\partial S^2} + rS \frac{\partial V}{\partial S} -rV &=0
\end{align*}

This is the famous Black-Scholes-Merton PDF, formally written with the
boundary conditions as follows

\begin{align*}
\frac{\partial c}{\partial t} + \frac{1}{2} \sigma^2 c^2 \frac{\partial^2 c}{\partial S^2} + rc \frac{\partial c}{\partial S} -rc &=0\\
c(0,t) &= 0\\ 
c(S_{+\infty}, t) &= S - Ke^{-r(T-t)}\\
c(S,T) & = max\{S-K,0\}
\end{align*}

This Black-Scholes-Merton PDE can be reduced to the heat equation using
the substitutions \(S = K e^x\),
\(t = T - \frac{\tau}{\frac{1}{2} \sigma^2}\), and
\(c(S, t) = K v(x, \tau)\). Let's derive the solution step by step in
full mathematical detail and show how this leads to the normal CDF.

\paragraph{Step 1: Substitutions}\label{step-1-substitutions}

We aim to reduce the BSM PDE: \[
\frac{\partial c}{\partial t} + \frac{1}{2} \sigma^2 S^2 \frac{\partial^2 c}{\partial S^2} + r S \frac{\partial c}{\partial S} - r c = 0
\]

to the heat equation. Using the substitutions:

\begin{itemize}
\tightlist
\item
  \(S = K e^x\), where \(x = \ln(S / K)\), and \(S \in (0, \infty)\)
  maps \(x \in (-\infty, \infty)\),
\item
  \(t = T - \frac{\tau}{\frac{1}{2} \sigma^2}\), so
  \(\tau = \frac{1}{2} \sigma^2 (T - t)\),
\item
  \(c(S, t) = K v(x, \tau)\), where \(v(x, \tau)\) is the transformed
  function.
\end{itemize}

\paragraph{Step 2: Derivative
Transformations}\label{step-2-derivative-transformations}

For \(c(S, t) = K v(x, \tau)\), we compute derivatives.

\begin{enumerate}
\def\labelenumi{\arabic{enumi}.}
\item
  The first derivative of \(c\) with respect to \(S\): \[
  \frac{\partial c}{\partial S} = \frac{\partial}{\partial S} \big(K v(x, \tau)\big) = K \frac{\partial v}{\partial x} \frac{\partial x}{\partial S},
  \] where \(x = \ln(S / K)\) implies
  \(\frac{\partial x}{\partial S} = \frac{1}{S}\). Thus: \[
  \frac{\partial c}{\partial S} = K \frac{\partial v}{\partial x} \frac{1}{S}.
  \]
\item
  The second derivative of \(c\) with respect to \(S\): \[
  \frac{\partial^2 c}{\partial S^2} = \frac{\partial}{\partial S} \left( K \frac{\partial v}{\partial x} \frac{1}{S} \right).
  \] Using the product rule: \[
  \frac{\partial^2 c}{\partial S^2} = K \frac{\partial^2 v}{\partial x^2} \frac{1}{S^2} - K \frac{\partial v}{\partial x} \frac{1}{S^2}.
  \]
\item
  The time derivative: \[
  \frac{\partial c}{\partial t} = K \frac{\partial v}{\partial \tau} \frac{\partial \tau}{\partial t}, \quad \text{and } \frac{\partial \tau}{\partial t} = -\frac{1}{\frac{1}{2} \sigma^2}.
  \]
\end{enumerate}

\paragraph{Step 3: Transforming the
PDE}\label{step-3-transforming-the-pde}

Substituting the above derivatives into the BSM PDE, we rewrite each
term.

\begin{enumerate}
\def\labelenumi{\arabic{enumi}.}
\item
  For \(\frac{\partial c}{\partial t}\): \[
  \frac{\partial c}{\partial t} = -\frac{1}{\frac{1}{2} \sigma^2} K \frac{\partial v}{\partial \tau}.
  \]
\item
  For \(\frac{\partial c}{\partial S}\): \[
  S \frac{\partial c}{\partial S} = S \cdot \left(K \frac{\partial v}{\partial x} \frac{1}{S}\right) = K \frac{\partial v}{\partial x}.
  \]
\item
  For \(\frac{\partial^2 c}{\partial S^2}\): \[
  \frac{1}{2} \sigma^2 S^2 \frac{\partial^2 c}{\partial S^2} = \frac{1}{2} \sigma^2 S^2 \left(K \frac{\partial^2 v}{\partial x^2} \frac{1}{S^2} - K \frac{\partial v}{\partial x} \frac{1}{S^2}\right) = \frac{1}{2} \sigma^2 K \frac{\partial^2 v}{\partial x^2}.
  \]
\end{enumerate}

Substituting all these into the BSM PDE: \[
-\frac{1}{\frac{1}{2} \sigma^2} K \frac{\partial v}{\partial \tau} + \frac{1}{2} \sigma^2 K \frac{\partial^2 v}{\partial x^2} + r K \frac{\partial v}{\partial x} - r K v = 0.
\]

Divide through by \(K\): \[
-\frac{\partial v}{\partial \tau} + \frac{\partial^2 v}{\partial x^2} + \frac{2r}{\sigma^2} \frac{\partial v}{\partial x} - \frac{2r}{\sigma^2} v = 0.
\]

To simplify, let \(v(x, \tau) = e^{\alpha x + \beta \tau} u(x, \tau)\),
where \(\alpha\) and \(\beta\) are constants. Substituting and choosing
\(\alpha = -\frac{r}{\sigma^2}\) and
\(\beta = -\frac{r^2}{2 \sigma^2}\), the equation reduces to: \[
\frac{\partial u}{\partial \tau} = \frac{\partial^2 u}{\partial x^2}.
\]

\paragraph{Step 4: Solving the Heat
Equation}\label{step-4-solving-the-heat-equation}

The heat equation
\(\frac{\partial u}{\partial \tau} = \frac{\partial^2 u}{\partial x^2}\)
has a well-known solution using Fourier methods: \[
u(x, \tau) = \frac{1}{\sqrt{2 \pi \tau}} \int_{-\infty}^\infty e^{-\frac{(x-y)^2}{2\tau}} f(y) \, dy,
\]

where \(f(y)\) is the initial condition.

For the BSM problem, the initial condition is the payoff: \[
f(y) = \max(e^y - 1, 0).
\]

Performing the integration leads to the final solution involving the
cumulative normal distribution function: \[
v(x, \tau) = N(d_1) - e^{-x} N(d_2),
\]

where: \[
d_1 = \frac{x + \frac{1}{2} \tau}{\sqrt{\tau}}, \quad d_2 = \frac{x - \frac{1}{2} \tau}{\sqrt{\tau}}.
\]

Transforming back to the original variables gives the Black-Scholes
formula: \[
C(S, t) = S e^{-q(T-t)} N(d_1) - K e^{-r(T-t)} N(d_2),
\] where: \[
d_1 = \frac{\ln(S / K) + (r - q + \frac{\sigma^2}{2})(T-t)}{\sigma \sqrt{T-t}}, \quad d_2 = d_1 - \sigma \sqrt{T-t}.
\]

Similarly, we can derive the price of a European put option:

\[
P = K e^{-rT} N(-d_2) - S e^{-qT} N(-d_1)
\]

Where: \[
d_1 = \frac{\ln(\frac{S}{K}) + (r - q + \frac{\sigma^2}{2})T}{\sigma \sqrt{T}}, \quad d_2 = d_1 - \sigma \sqrt{T}
\]

\subsubsection{Asymptotic Behavior of the BSM formula for call and put
options}\label{asymptotic-behavior-of-the-bsm-formula-for-call-and-put-options}

What if \(K\rightarrow 0\)? In that case,

\begin{enumerate}
\def\labelenumi{\arabic{enumi}.}
\tightlist
\item
  \(\ln(S_0/K)\rightarrow \infty\), causing \(d_1 \rightarrow \infty\)
  and \(d_2 \rightarrow \infty\)\\
\item
  The cdf \(N(d_1)\rightarrow 1\) and \(N(d_2)\rightarrow 1\)\\
\item
  The second term \(Ke^{-rT}N(d_2)\rightarrow 0\) as \(K\rightarrow 0\)
\end{enumerate}

In this case, the price of a call option \(C\rightarrow S_0\) and the
price of a put option \(P \rightarrow 0\)

\subsection{Greeks: Delta and Gamma}\label{greeks-delta-and-gamma}

\textbf{Delta} (\(\Delta\)) is the sensitivity of the option price to
changes in the underlying asset price:

\[
\Delta = \frac{\partial C}{\partial S}\approx \frac{C(S_0 + h) - C(S_0 - h)}{2h}
\]

This is the \textbf{central difference approximation}, which provides a
more accurate estimate of delta compared to the forward or backward
difference methods.

\begin{itemize}
\tightlist
\item
  \(C(S_0 + h)\): Calculate the option price with the spot price
  increased by \(h\).
\item
  \(C(S_0 - h)\): Calculate the option price with the spot price
  decreased by \(h\).
\end{itemize}

\textbf{Gamma} (\(\Gamma\)) measures the rate of change of delta with
respect to the underlying asset price:

\[
\Gamma = \frac{\partial^2 C}{\partial S^2}\approx \frac{\Delta(S_0 + h) - \Delta(S_0 - h)}{2h}\approx \frac{C(S_0 + h) - 2C(S_0) + C(S_0 - h)}{h^2}
\]

Gamma (\(\Gamma\)) measures the rate of change of delta (\(\Delta\))
with respect to the underlying spot price (\(S_0\)).

\begin{itemize}
\tightlist
\item
  \(C(S_0 + h)\): Option price with the spot price increased by \(h\).
\item
  \(C(S_0)\): Option price at the current spot price.
\item
  \(C(S_0 - h)\): Option price with the spot price decreased by \(h\).
\end{itemize}

\textbf{Relationship Between Delta and Gamma:}

\begin{itemize}
\tightlist
\item
  Gamma represents how much delta changes for a small change in \(S_0\).
\item
  If gamma is high, delta is more sensitive to changes in \(S_0\), which
  is important for hedging strategies.
\end{itemize}

\subsection{Implementation}\label{implementation}

\subsubsection{Notation}\label{notation}

\begin{itemize}
\tightlist
\item
  \(S\): Spot price of the stock.
\item
  \(K\): Strike price of the option.
\item
  \(T\): Time to maturity (in years).
\item
  \(r\): Risk-free rate (continuously compounded).
\item
  \(q\): Dividend yield (continuously compounded).
\item
  \(\sigma\): Volatility of the stock.
\item
  \(N(\cdot)\): Cumulative distribution function of the standard normal
  distribution.
\end{itemize}

\begin{Shaded}
\begin{Highlighting}[]
\ImportTok{from}\NormalTok{ dataclasses }\ImportTok{import}\NormalTok{ dataclass}
\ImportTok{import}\NormalTok{ numpy }\ImportTok{as}\NormalTok{ np}
\ImportTok{from}\NormalTok{ scipy.stats }\ImportTok{import}\NormalTok{ norm}

\AttributeTok{@dataclass}
\KeywordTok{class}\NormalTok{ Equity:}
\NormalTok{    spot: }\BuiltInTok{float}
\NormalTok{    dividend\_yield: }\BuiltInTok{float}
\NormalTok{    volatility: }\BuiltInTok{float}

\AttributeTok{@dataclass}
\KeywordTok{class}\NormalTok{ EquityOption:}
\NormalTok{    strike: }\BuiltInTok{float}
\NormalTok{    time\_to\_maturity: }\BuiltInTok{float}
\NormalTok{    put\_call: }\BuiltInTok{str}

\AttributeTok{@dataclass}
\KeywordTok{class}\NormalTok{ EquityForward:}
\NormalTok{    strike: }\BuiltInTok{float}
\NormalTok{    time\_to\_maturity: }\BuiltInTok{float}

\KeywordTok{def}\NormalTok{ bsm(underlying: Equity, option: EquityOption, rate: }\BuiltInTok{float}\NormalTok{) }\OperatorTok{{-}\textgreater{}} \BuiltInTok{float}\NormalTok{:}
\NormalTok{    S }\OperatorTok{=}\NormalTok{ underlying.spot}
\NormalTok{    K }\OperatorTok{=}\NormalTok{ option.strike}
\NormalTok{    T }\OperatorTok{=}\NormalTok{ option.time\_to\_maturity}
\NormalTok{    r }\OperatorTok{=}\NormalTok{ rate}
\NormalTok{    q }\OperatorTok{=}\NormalTok{ underlying.dividend\_yield}
\NormalTok{    sigma }\OperatorTok{=}\NormalTok{ underlying.volatility}

    \CommentTok{\# Handle edge case where strike is effectively zero}
    \ControlFlowTok{if}\NormalTok{ K }\OperatorTok{\textless{}} \FloatTok{1e{-}8}\NormalTok{:}
        \ControlFlowTok{if}\NormalTok{ option.put\_call.lower() }\OperatorTok{==} \StringTok{"call"}\NormalTok{:}
            \ControlFlowTok{return}\NormalTok{ S }
        \ControlFlowTok{else}\NormalTok{:}
            \ControlFlowTok{return} \FloatTok{0.0}

\NormalTok{    d1 }\OperatorTok{=}\NormalTok{ (np.log(S }\OperatorTok{/}\NormalTok{ K) }\OperatorTok{+}\NormalTok{ (r }\OperatorTok{{-}}\NormalTok{ q }\OperatorTok{+} \FloatTok{0.5} \OperatorTok{*}\NormalTok{ sigma}\OperatorTok{**}\DecValTok{2}\NormalTok{) }\OperatorTok{*}\NormalTok{ T) }\OperatorTok{/}\NormalTok{ (sigma }\OperatorTok{*}\NormalTok{ np.sqrt(T))}
\NormalTok{    d2 }\OperatorTok{=}\NormalTok{ d1 }\OperatorTok{{-}}\NormalTok{ sigma }\OperatorTok{*}\NormalTok{ np.sqrt(T)}

    \ControlFlowTok{if}\NormalTok{ option.put\_call.lower() }\OperatorTok{==} \StringTok{"call"}\NormalTok{:}
\NormalTok{        price }\OperatorTok{=}\NormalTok{ S }\OperatorTok{*}\NormalTok{ np.exp(}\OperatorTok{{-}}\NormalTok{q }\OperatorTok{*}\NormalTok{ T) }\OperatorTok{*}\NormalTok{ norm.cdf(d1) }\OperatorTok{\textbackslash{}}
                \OperatorTok{{-}}\NormalTok{ K }\OperatorTok{*}\NormalTok{ np.exp(}\OperatorTok{{-}}\NormalTok{r }\OperatorTok{*}\NormalTok{ T) }\OperatorTok{*}\NormalTok{ norm.cdf(d2)}
    \ControlFlowTok{elif}\NormalTok{ option.put\_call.lower() }\OperatorTok{==} \StringTok{"put"}\NormalTok{:}
\NormalTok{        price }\OperatorTok{=}\NormalTok{ K }\OperatorTok{*}\NormalTok{ np.exp(}\OperatorTok{{-}}\NormalTok{r }\OperatorTok{*}\NormalTok{ T) }\OperatorTok{*}\NormalTok{ norm.cdf(}\OperatorTok{{-}}\NormalTok{d2) }\OperatorTok{\textbackslash{}}
                \OperatorTok{{-}}\NormalTok{ S }\OperatorTok{*}\NormalTok{ np.exp(}\OperatorTok{{-}}\NormalTok{q }\OperatorTok{*}\NormalTok{ T) }\OperatorTok{*}\NormalTok{ norm.cdf(}\OperatorTok{{-}}\NormalTok{d1)}
    \ControlFlowTok{else}\NormalTok{:}
        \ControlFlowTok{raise} \PreprocessorTok{ValueError}\NormalTok{(}\StringTok{"Invalid option type. Must be \textquotesingle{}call\textquotesingle{} or \textquotesingle{}put\textquotesingle{}."}\NormalTok{)}

    \ControlFlowTok{return}\NormalTok{ price}

\KeywordTok{def}\NormalTok{ delta(underlying: Equity, option: EquityOption, rate: }\BuiltInTok{float}\NormalTok{) }\OperatorTok{{-}\textgreater{}} \BuiltInTok{float}\NormalTok{:}
\NormalTok{    bump }\OperatorTok{=} \FloatTok{0.01} \OperatorTok{*}\NormalTok{ underlying.spot}
\NormalTok{    bumped\_up }\OperatorTok{=}\NormalTok{ Equity(spot}\OperatorTok{=}\NormalTok{underlying.spot }\OperatorTok{+}\NormalTok{ bump, }
\NormalTok{                       dividend\_yield}\OperatorTok{=}\NormalTok{underlying.dividend\_yield, }
\NormalTok{                       volatility}\OperatorTok{=}\NormalTok{underlying.volatility)}
\NormalTok{    bumped\_down }\OperatorTok{=}\NormalTok{ Equity(spot}\OperatorTok{=}\NormalTok{underlying.spot }\OperatorTok{{-}}\NormalTok{ bump, }
\NormalTok{                         dividend\_yield}\OperatorTok{=}\NormalTok{underlying.dividend\_yield, }
\NormalTok{                         volatility}\OperatorTok{=}\NormalTok{underlying.volatility)}
\NormalTok{    price\_up }\OperatorTok{=}\NormalTok{ bsm(bumped\_up, option, rate)}
\NormalTok{    price\_down }\OperatorTok{=}\NormalTok{ bsm(bumped\_down, option, rate)}
    \ControlFlowTok{return}\NormalTok{ (price\_up }\OperatorTok{{-}}\NormalTok{ price\_down) }\OperatorTok{/}\NormalTok{ (}\DecValTok{2} \OperatorTok{*}\NormalTok{ bump)}

\KeywordTok{def}\NormalTok{ gamma(underlying: Equity, option: EquityOption, rate: }\BuiltInTok{float}\NormalTok{) }\OperatorTok{{-}\textgreater{}} \BuiltInTok{float}\NormalTok{:}
\NormalTok{    bump }\OperatorTok{=} \FloatTok{0.01} \OperatorTok{*}\NormalTok{ underlying.spot}
\NormalTok{    bumped\_up }\OperatorTok{=}\NormalTok{ Equity(spot}\OperatorTok{=}\NormalTok{underlying.spot }\OperatorTok{+}\NormalTok{ bump, }
\NormalTok{                       dividend\_yield}\OperatorTok{=}\NormalTok{underlying.dividend\_yield, }
\NormalTok{                       volatility}\OperatorTok{=}\NormalTok{underlying.volatility)}
\NormalTok{    bumped\_down }\OperatorTok{=}\NormalTok{ Equity(spot}\OperatorTok{=}\NormalTok{underlying.spot }\OperatorTok{{-}}\NormalTok{ bump, }
\NormalTok{                         dividend\_yield}\OperatorTok{=}\NormalTok{underlying.dividend\_yield, }
\NormalTok{                         volatility}\OperatorTok{=}\NormalTok{underlying.volatility)}
\NormalTok{    original\_price }\OperatorTok{=}\NormalTok{ bsm(underlying, option, rate)}
\NormalTok{    price\_up }\OperatorTok{=}\NormalTok{ bsm(bumped\_up, option, rate)}
\NormalTok{    price\_down }\OperatorTok{=}\NormalTok{ bsm(bumped\_down, option, rate)}
    \ControlFlowTok{return}\NormalTok{ (price\_up }\OperatorTok{{-}} \DecValTok{2} \OperatorTok{*}\NormalTok{ original\_price }\OperatorTok{+}\NormalTok{ price\_down) }\OperatorTok{/}\NormalTok{ (bump}\OperatorTok{**}\DecValTok{2}\NormalTok{)}

\KeywordTok{def}\NormalTok{ fwd(underlying: Equity, forward: EquityForward, rate: }\BuiltInTok{float}\NormalTok{) }\OperatorTok{{-}\textgreater{}} \BuiltInTok{float}\NormalTok{:}
\NormalTok{    S }\OperatorTok{=}\NormalTok{ underlying.spot}
\NormalTok{    K }\OperatorTok{=}\NormalTok{ forward.strike}
\NormalTok{    T }\OperatorTok{=}\NormalTok{ forward.time\_to\_maturity}
\NormalTok{    r }\OperatorTok{=}\NormalTok{ rate}
\NormalTok{    q }\OperatorTok{=}\NormalTok{ underlying.dividend\_yield}
\NormalTok{    forward\_price }\OperatorTok{=}\NormalTok{ S }\OperatorTok{*}\NormalTok{ np.exp((r }\OperatorTok{{-}}\NormalTok{ q) }\OperatorTok{*}\NormalTok{ T) }\OperatorTok{{-}}\NormalTok{ K}

    \ControlFlowTok{return}\NormalTok{ forward\_price}

\KeywordTok{def}\NormalTok{ check\_put\_call\_parity(}
\NormalTok{    underlying: Equity, }
\NormalTok{    call\_option: EquityOption, }
\NormalTok{    put\_option: EquityOption, }
\NormalTok{    rate: }\BuiltInTok{float}
\NormalTok{    ) }\OperatorTok{{-}\textgreater{}} \BuiltInTok{bool}\NormalTok{:}

\NormalTok{    call\_price }\OperatorTok{=}\NormalTok{ bsm(underlying, call\_option, rate)}
\NormalTok{    put\_price }\OperatorTok{=}\NormalTok{ bsm(underlying, put\_option, rate)}
\NormalTok{    S }\OperatorTok{=}\NormalTok{ underlying.spot}
\NormalTok{    K }\OperatorTok{=}\NormalTok{ call\_option.strike}
\NormalTok{    T }\OperatorTok{=}\NormalTok{ call\_option.time\_to\_maturity}
\NormalTok{    r }\OperatorTok{=}\NormalTok{ rate}
\NormalTok{    q }\OperatorTok{=}\NormalTok{ underlying.dividend\_yield}

\NormalTok{    parity\_lhs }\OperatorTok{=}\NormalTok{ call\_price }\OperatorTok{{-}}\NormalTok{ put\_price}
\NormalTok{    parity\_rhs }\OperatorTok{=}\NormalTok{ S }\OperatorTok{*}\NormalTok{ np.exp(}\OperatorTok{{-}}\NormalTok{q }\OperatorTok{*}\NormalTok{ T) }\OperatorTok{{-}}\NormalTok{ K }\OperatorTok{*}\NormalTok{ np.exp(}\OperatorTok{{-}}\NormalTok{r }\OperatorTok{*}\NormalTok{ T)}

    \ControlFlowTok{return}\NormalTok{ np.isclose(parity\_lhs, parity\_rhs, atol}\OperatorTok{=}\FloatTok{1e{-}4}\NormalTok{)}
\end{Highlighting}
\end{Shaded}

\subsubsection{Example Usage}\label{example-usage}

Say, we want to price a call option on an equity with spot price
\(S_0 = 450\) with dividend yield \(q=1.4\%\), and volatility \(14\%\).
The strike price of the call is \(K=470\), with time to maturity in
years \(T=0.23\) and the risk free rate \(r = 0.05\). Next, we want to
see the asymptotic behavior of the call option if the strike price
\(K\rightarrow 0\) with interest rate 0. Next, we want to price a put
option on the same equity but strike price \(K=500\), time to maturity
in years \(T=0.26\) and interest rate is 0. Finally, we want to check if
the put-call parity relationship is hold. In each case, we consider
\(h=0.01\) a bump or small change in the stock price.

\begin{Shaded}
\begin{Highlighting}[]
\ControlFlowTok{if} \VariableTok{\_\_name\_\_} \OperatorTok{==} \StringTok{"\_\_main\_\_"}\NormalTok{:}
\NormalTok{    eq }\OperatorTok{=}\NormalTok{ Equity(}\DecValTok{450}\NormalTok{, }\FloatTok{0.014}\NormalTok{, }\FloatTok{0.14}\NormalTok{)}
\NormalTok{    option\_call }\OperatorTok{=}\NormalTok{ EquityOption(}\DecValTok{470}\NormalTok{, }\FloatTok{0.23}\NormalTok{, }\StringTok{"call"}\NormalTok{)}
\NormalTok{    option\_put }\OperatorTok{=}\NormalTok{ EquityOption(}\DecValTok{500}\NormalTok{, }\FloatTok{0.26}\NormalTok{, }\StringTok{"put"}\NormalTok{)}
    
    \BuiltInTok{print}\NormalTok{(bsm(eq, option\_call, }\FloatTok{0.05}\NormalTok{))  }
    \BuiltInTok{print}\NormalTok{(bsm(eq, EquityOption(}\FloatTok{1e{-}15}\NormalTok{, }\FloatTok{0.26}\NormalTok{, }\StringTok{"call"}\NormalTok{), }\FloatTok{0.0}\NormalTok{))  }
    \BuiltInTok{print}\NormalTok{(bsm(Equity(}\DecValTok{450}\NormalTok{, }\FloatTok{0.0}\NormalTok{, }\FloatTok{1e{-}9}\NormalTok{), option\_put, }\FloatTok{0.0}\NormalTok{))  }

    \CommentTok{\# Check put{-}call parity}
\NormalTok{    eq }\OperatorTok{=}\NormalTok{ Equity(}\DecValTok{450}\NormalTok{, }\FloatTok{0.015}\NormalTok{, }\FloatTok{0.15}\NormalTok{)}
\NormalTok{    option\_call }\OperatorTok{=}\NormalTok{ EquityOption(}\DecValTok{470}\NormalTok{, }\FloatTok{0.26}\NormalTok{, }\StringTok{"call"}\NormalTok{)}
\NormalTok{    option\_put }\OperatorTok{=}\NormalTok{ EquityOption(}\DecValTok{470}\NormalTok{, }\FloatTok{0.26}\NormalTok{, }\StringTok{"put"}\NormalTok{)}
    \BuiltInTok{print}\NormalTok{(check\_put\_call\_parity(eq, option\_call, option\_put, }\FloatTok{0.05}\NormalTok{)) }
\end{Highlighting}
\end{Shaded}

\begin{verbatim}
5.834035584709994
450
50.0
True
\end{verbatim}

\subsection{References}\label{references}

\begin{itemize}
\tightlist
\item
  Karatzas, I., \& Shreve, S. E. (1991). \emph{Brownian Motion and
  Stochastic Calculus}.\\
\item
  Options, Futures, and Other Derivatives by John C. Hull\\
\item
  Arbitrage Theory in Continuous Time Book by Tomas Björk
\end{itemize}

\textbf{Share on}

\faIcon{facebook}

\faIcon{linkedin}

\faIcon{twitter}

\protect\phantomsection\label{fb-root}

\textbf{You may also like}


\printbibliography



\end{document}
